%! Dimensions of the Four Subspaces — Unit 3, Lesson 3.5 — Homework (Answer Key)
\documentclass[10pt]{article}
\usepackage{linalg-article}
\usepackage{linalg-key}   % pulls in -boxes; do NOT also load -boxes
\newcommand{\TallMath}[1]{$\displaystyle #1\rule[-1.4em]{0pt}{3.2em}$}
\newcommand{\xx}{\mathbf{x}}
\newcommand{\yy}{\mathbf{y}}
\newcommand{\svec}{\mathbf{s}}
\newcommand{\zero}{\mathbf{0}}
\newcommand{\T}{\mathsf{T}}

\begin{document}

\pageheader{Unit 3, Lesson 3.5}{Homework --- Answer Key}
\namedateperiod

\vspace{0.12in}

\begin{notesbox}{Practice}
Remember the four subspaces: in $\mathbb{R}^n$ the row space $C(A^{\T})$ ($\dim r$) and nullspace $N(A)$
($\dim n-r$); in $\mathbb{R}^m$ the column space $C(A)$ ($\dim r$) and left nullspace $N(A^{\T})$
($\dim m-r$). Always \emph{justify} and give each dimension.

\begin{enumerate}[leftmargin=1.9em, itemsep=11pt]
  \item \textbf{Dimensions from $r$.} Fill in the four dimensions for each matrix.
    \begin{center}\small
    \renewcommand{\arraystretch}{1.5}
    \begin{tabular}{@{}l c c c c@{}}
    \textbf{Matrix} & $\dim C(A)$ & $\dim C(A^{\T})$ & $\dim N(A)$ & $\dim N(A^{\T})$ \\
    \hline
    (a) $3\times3$, $r=3$ & \ans{$3$} & \ans{$3$} & \ans{$0$} & \ans{$0$} \\
    (b) $2\times4$, $r=2$ & \ans{$2$} & \ans{$2$} & \ans{$2$} & \ans{$0$} \\
    (c) $5\times3$, $r=2$ & \ans{$2$} & \ans{$2$} & \ans{$1$} & \ans{$3$} \\
    \end{tabular}
    \end{center}
  \item \textbf{All four from one reduction.} For $A=\begin{bmatrix} 1 & 2 & 1 \\ 1 & 3 & 2 \\ 2 & 5 & 3 \end{bmatrix}$
        (row~3 $=$ row~1 $+$ row~2), reduce, then give a basis \emph{and} dimension for $C(A)$, $C(A^{\T})$,
        $N(A)$, and $N(A^{\T})$. Check $r+(n-r)=n$ and $r+(m-r)=m$.
        \ansline{Reduce to $\left[\begin{smallmatrix} 1&0&-1\\0&1&1\\0&0&0 \end{smallmatrix}\right]$: $m=n=3$, $r=2$. \textbf{$C(A)\subseteq\mathbb{R}^3$:} pivot columns $\begin{bsmallmatrix} 1\\1\\2 \end{bsmallmatrix},\begin{bsmallmatrix} 2\\3\\5 \end{bsmallmatrix}$, $\dim=2$. \textbf{$C(A^{\T})\subseteq\mathbb{R}^3$:} rows $(1,0,-1),(0,1,1)$, $\dim=2$. \textbf{$N(A)\subseteq\mathbb{R}^3$:} $\svec=\begin{bsmallmatrix} 1\\-1\\1 \end{bsmallmatrix}$, $\dim=n-r=1$. \textbf{$N(A^{\T})\subseteq\mathbb{R}^3$:} $\yy=\begin{bsmallmatrix} 1\\1\\-1 \end{bsmallmatrix}$, $\dim=m-r=1$. Checks: $2+1=3=n$, $2+1=3=m$.}
  \item \textbf{Read the redundancies.} For the matrix in problem~2, the nullspace vector is
        $\svec=\begin{bsmallmatrix} 1\\-1\\1 \end{bsmallmatrix}$ and a left-nullspace vector is
        $\yy=\begin{bsmallmatrix} 1\\1\\-1 \end{bsmallmatrix}$. Write the \emph{column} dependency $\svec$
        names and the \emph{row} dependency $\yy$ names. \ansline{$A\svec=\zero$: $\mathbf{a}_1-\mathbf{a}_2+\mathbf{a}_3=\zero$, so column~3 $=$ column~2 $-$ column~1 (i.e.\ $\begin{bsmallmatrix} 1\\2\\3 \end{bsmallmatrix}=\begin{bsmallmatrix} 2\\3\\5 \end{bsmallmatrix}-\begin{bsmallmatrix} 1\\1\\2 \end{bsmallmatrix}$). $A^{\T}\yy=\zero$: row$_1+$row$_2-$row$_3=\zero$, so row~3 $=$ row~1 $+$ row~2.}
  \item \textbf{Explain.} A matrix has ``full column rank'' when $r=n$ and ``full row rank'' when $r=m$. Using
        the dimension formulas, explain why full column rank forces $N(A)=\{\zero\}$ and full row rank forces
        $N(A^{\T})=\{\zero\}$. \ansline{$\dim N(A)=n-r$; full column rank $r=n$ makes this $0$, so $N(A)=\{\zero\}$ (no free columns --- the columns are independent). $\dim N(A^{\T})=m-r$; full row rank $r=m$ makes this $0$, so $N(A^{\T})=\{\zero\}$ (no redundant rows --- the rows are independent).}
\end{enumerate}
\end{notesbox}

\begin{extensionbox}
\textbf{One number, four dimensions.}
\begin{enumerate}[label=(\alph*), leftmargin=1.8em, itemsep=5pt]
  \item \textbf{Why row rank $=$ column rank.} In one or two sentences, explain why the $r$ pivots of a
        reduction count both the independent columns and the independent rows.
  \item \textbf{An invertible matrix.} Let $A$ be an invertible $n\times n$ matrix ($r=n$). State all four
        subspaces and their dimensions. What happens to the nullspace and the left nullspace?
\end{enumerate}
\ansline{(a) A reduction has $r$ pivots, one per pivot column and one per nonzero row. The $r$ pivot columns are independent and span $C(A)$; the $r$ nonzero reduced rows are independent and span the row space. The same $r$ pivots count both, so $\dim C(A)=\dim C(A^{\T})=r$.}
\ansline{(b) With $r=n=m$: $C(A)=\mathbb{R}^n$ and $C(A^{\T})=\mathbb{R}^n$, each $\dim=n$; $N(A)=\{\zero\}$ and $N(A^{\T})=\{\zero\}$, each $\dim=0$. Both nullspaces collapse to the zero vector --- an invertible matrix has no redundant columns and no redundant rows.}
\end{extensionbox}

\begin{spiralbox}
\textbf{Coming next --- Lesson 3.6: The Fundamental Theorem of Linear Algebra.} You now have all four
subspaces and their dimensions ($r$, $r$, $n-r$, $m-r$) from one number $r$. Next we assemble them into the
``big picture'' of a matrix --- two subspaces in each room --- and see how the row space and nullspace (and
the column space and left nullspace) sit at \emph{right angles}, opening the door to orthogonality in Unit~4.
\end{spiralbox}

\begin{teachernote}
Problem~1 is pure counting --- the fastest check that all four formulas ($r,r,n-r,m-r$) are automatic;
watch for $m$/$n$ swaps in (c) ($5\times3$: $\dim N(A^{\T})=m-r=3$, not $1$). Problem~2 is the core skill,
all four bases from one reduction. Problem~3 makes both nullspaces concrete: a column dependency from
$\svec$, a row dependency from $\yy$. Extension~(b) is the capstone case --- invertible means both nullspaces
are $\{\zero\}$ and both spaces fill $\mathbb{R}^n$; it sets up 3.6's big picture.
\end{teachernote}

\end{document}
